\documentclass[11pt]{article}
\usepackage[margin=1in]{geometry}
\usepackage{amsmath,amssymb,amsthm,mathtools}
\usepackage{microtype}
\usepackage{enumitem}
\usepackage{booktabs}
\usepackage[hidelinks]{hyperref}
\setlist{nosep}
\newtheorem{theorem}{Theorem}[section]
\newtheorem{proposition}[theorem]{Proposition}
\newtheorem{corollary}[theorem]{Corollary}
\newtheorem{lemma}[theorem]{Lemma}
\theoremstyle{definition}
\newtheorem{definition}[theorem]{Definition}
\newtheorem{warning}[theorem]{Warning}
\theoremstyle{remark}
\newtheorem{remark}[theorem]{Remark}
\newcommand{\Ext}{\operatorname{Ext}}
\newcommand{\osc}{\operatorname{osc}}
\newcommand{\op}{\mathrm{op}}
\newcommand{\preceqL}{\preceq}
\title{Proposed Manuscript Insertions III\\
End Rescaling, Adaptive Plumbing Response, and Finite-Horizon Transfer}
\author{Research draft for the marked-width MLC program}
\date{July 24, 2026}

\begin{document}
\maketitle

\begin{abstract}
We construct a canonical direct-limit rescaling surface for every infinite
address of the first direct-decoration pullback system and prove that every
finite, shift-consistent diagram on these surfaces has an exact realization by
finite polynomial cylinders.  We then use the source ray/equipotential
quadrilaterals to obtain linear serial resistance growth and an explicit
operator-norm estimate comparing a deep physical response matrix with its
nodal normalization.  Any preassigned additive plumbing error can be achieved
by a finite depth choice.  Consequently, a source-licensed strict root
transition transfers through a direct-decoration virtual end with the same
homogeneous contraction coefficient.  These results do not identify the
normalized fibers with Dudko's proposed actual planar Cantor family and do not
prove MLC.
\end{abstract}

\section{Direct-decoration data}
Let $P$ be a closed Jordan puzzle disk, let $g=f^q$, and let
$\{Q_a:a\in\mathcal A\}$ be finitely many closed Jordan pullbacks such that
$g:Q_a\to P$ is conformal.  Write $\phi_a:P\to Q_a$ for the inverse branch.
For $w=a_1\cdots a_m$, put
\[
 \phi_w=\phi_{a_1}\circ\cdots\circ\phi_{a_m},\quad
 Q_w=\phi_w(P),\quad h_w=g^m|_{Q_w}.
\]
Let
\[
 F=f^r|_A:A\to B,\qquad K\Subset A\Subset B\Subset\operatorname{int}P
\]
be a connected quadratic-like restriction, and let $\jmath:A\hookrightarrow B$
be inclusion.  For a finite word $w$, define
\[
 A_w=h_w^{-1}(A),\quad K_w=h_w^{-1}(K),
\]
\[
 G_w=F\circ h_w=f^{q|w|+r}|_{A_w},\qquad
 I_w=\jmath\circ h_w=f^{q|w|}|_{A_w}.
\]
The relative map is the root map:
\[
 G_w\circ I_w^{-1}=F\circ\jmath^{-1}.
\]
Thus the common univalent prefix does not alter the relative critical word.

\section{Canonical end rescaling}
Fix $\omega=a_1a_2\cdots\in\mathcal A^{\mathbb N}$ and write
$w_m=\omega|m$.  Define the geometric-prefix map
\[
 t_m=h_{w_{m+1}}^{-1}\circ h_{w_m}:A_{w_m}\to A_{w_{m+1}}.
\]
Then
\[
 h_{w_{m+1}}t_m=h_{w_m},\qquad
 G_{w_{m+1}}t_m=G_{w_m},\qquad
 I_{w_{m+1}}t_m=I_{w_m}.
\]

\begin{theorem}[Canonical end-rescaling surface]
The direct limit
\[
 \widehat A_\omega=\varinjlim(A_{w_m},t_m)
\]
is a Riemann surface equipped with a biholomorphism
\[
 h_\omega:\widehat A_\omega\to A.
\]
The finite pairs descend to
\[
 \widehat G_\omega=F\circ h_\omega,
 \qquad
 \widehat I_\omega=\jmath\circ h_\omega.
\]
The direct-limit carrier $\widehat K_\omega$ is mapped homeomorphically to
$K$.  Every finite marked datum transported by the maps $h_{w_m}$ is exactly
the corresponding root datum.
\end{theorem}

\begin{proof}
Take the quotient of $\bigsqcup_m A_{w_m}$ by the equivalence generated by
$z\sim t_m(z)$.  Since every transition is conformal, the quotient has the
unique Riemann-surface structure making the level maps conformal.  The formula
$h_\omega([z,m])=h_{w_m}(z)$ is well-defined and gives a conformal bijection
onto $A$.  Compatibility of the finite pairs and markings gives the remaining
claims.
\end{proof}

\begin{remark}[Scope]
The theorem constructs a normalized rescaling pro-object from actual finite
pullbacks.  It does not identify the planar impression
$E_\omega=\bigcap_mQ_{\omega|m}$ with a nondegenerate small Julia continuum.
\end{remark}

\subsection{Dynamic shift}
At finite levels,
\[
 g:A_{a_1\cdots a_{m+1}}\to A_{a_2\cdots a_{m+1}},\qquad
 h_{a_2\cdots a_{m+1}}g=h_{a_1\cdots a_{m+1}}.
\]

\begin{theorem}[End-shift naturality]
The finite maps induce a biholomorphism
\[
 \widehat S_\omega:\widehat A_\omega\to\widehat A_{\sigma\omega}
\]
with
\[
 h_{\sigma\omega}\widehat S_\omega=h_\omega.
\]
Hence the dynamic shift is the identity in normalized coordinates and commutes
with the descended covering/immersion pair and all transported finite markings.
\end{theorem}

\begin{warning}
The dynamic shift is induced by the actual iterate $f^q$.  The geometric-prefix
map $t_m$ is a conformal presentation change.  They become equal only after
normalization and must remain separately typed in the proof ledger.
\end{warning}

\section{Finite physical realization}
A finite normalized diagram is called \emph{depth graded} if its nodes carry
integers $s(v)$ such that root, presentation, and enlargement edges preserve
$s$, while a forward shift edge $u\to v$ satisfies $s(u)=s(v)+1$.

\begin{theorem}[Exact realization of a depth-graded diagram]
Let $\mathcal D$ be a finite depth-graded normalized diagram whose root
restrictions and marked squares are actual.  For all sufficiently large $M$,
represent node $v$ at physical depth $m(v)=M+s(v)$.  Then:
\begin{enumerate}
\item root branches are represented by the retimed pairs
$(f^{qm(v)+r_e},f^{qm(v)})$;
\item shift edges are represented by the actual map $f^q$;
\item every purely dynamical path is a genuine polynomial iterate on an actual
nonempty domain, with the exact chronological critical word;
\item presentation changes remain explicit conformal coordinate maps;
\item all declared component, sheet, obstacle, localization, prime-end,
relative-cover, response, Value, and enlargement data commute exactly.
\end{enumerate}
\end{theorem}

\begin{proof}
Pull every root object back by the cylinder chart at depth $m(v)$.  The grading
makes each shift source exactly one level deeper than its target.  The formulas
for $t_m$ and the finite shift give all generating squares.  The critical-spine
telescope and marked finite-horizon assembly apply to the dynamical
compositions.  Finiteness permits one common $M$.
\end{proof}

\begin{remark}
If two paths into one node have different net shift counts, the diagram must be
unfolded by duplicating that node.  Acyclicity alone does not imply physical
depth consistency.
\end{remark}

\section{Intrinsic finite pseudo-puzzle atlas}
Let $H:U\to V$ be an externally marked connected quadratic-like map.  Fix a
finite source-authorized list of rational external accesses, equipotential
cross-sections, and pullback depths.  Construct the finite access/equipotential
graph in the prime-end external class and take its actual pullbacks under $H$.

\begin{theorem}[Intrinsic finite atlas]
The resulting finite pseudo-puzzle consists of actual component domains and
actual restrictions of iterates $H^n$.  It retains exact return times, total
and local degrees, chronological critical words, component and sheet routes,
obstacles, prime-end phase, and cyclic order.  It is functorial under an
externally marked conjugacy of the declared finite atlas.
\end{theorem}

\begin{proof}
The external marking fixes the phase of the degree-two external covering.  The
finite graph is therefore defined in the external prime-end class.  Pullbacks,
components, and return exclusions are read from actual iterates of $H$.
No uniqueness of an arbitrary hybrid conjugacy is used.
\end{proof}

\begin{corollary}[Finite recursive compiler]
Every finite source-supplied decorated chain of externally marked connected
quadratic-like maps admits a prefix-local recursive realization on the
normalized end surfaces.  Every finite multi-scale diagram has an exact finite
physical realization.  The construction is covariant under marked changes of
presentation and does not assume rigidity.
\end{corollary}

\section{Exact fiberwise geometry}
Let $X\subset A\setminus K$ be a marked complement localization, let $\Gamma$
be a declared curve family, and transport all curves and electrodes by $h_w$
and $h_\omega$.

\begin{theorem}[Exact width and response]
For every finite word and infinite address,
\[
 \Ext_{X_w}(h_w^{-1}\Gamma)=\Ext_X(\Gamma)
 =\Ext_{\widehat X_\omega}(h_\omega^{-1}\Gamma).
\]
If $\Lambda_X$ is the finite-electrode Dirichlet-to-Neumann matrix, then
\[
 \Lambda_{X_w}=\Lambda_X=\Lambda_{\widehat X_\omega}.
\]
Parallel gluing and Schur-complement elimination commute with these equalities.
\end{theorem}

\begin{proof}
Extremal length and two-dimensional Dirichlet energy are conformal invariants.
The marked maps preserve electrode order and orientation.  The response gluing
operations are algebraic.
\end{proof}

\begin{corollary}
Any finite WAD, response, or Value state whose entire defining incidence,
frames, coefficient rows, and future observations are transported by the
marked conformal map is an identity state along the normalized address tail.
\end{corollary}

\section{Serial ray/equipotential necks}
For each letter $a$, choose a source-authorized quadrilateral
$C_a\subset P\setminus\operatorname{int}Q_a$ whose horizontal sides are
parent and child equipotential cross-sections and whose vertical sides are the
two retained ray/access walls.  Let $\mathcal V_a$ join the horizontal sides
and put
\[
 R_a=\Ext_{C_a}(\mathcal V_a)>0,\qquad
 R_*=\min_{a\in\mathcal A}R_a>0.
\]
For $w=a_1\cdots a_m$, set
\[
 C_{w,j}=\phi_{a_1\cdots a_{j-1}}(C_{a_j}).
\]
Call the localization serially exhaustive if every authorized deep-to-parent
curve contains a crossing of every $C_{w,j}$ in order.

\begin{theorem}[Serial resistance]
Under serial exhaustion, the effective resistance and conductance satisfy
\[
 R(w)\ge\sum_{j=1}^mR_{a_j}\ge mR_*,
\]
\[
 c(w)\le\left(\sum_{j=1}^mR_{a_j}\right)^{-1}
 \le\frac1{mR_*}.
\]
\end{theorem}

\begin{proof}
The pulled cells have disjoint interiors, are conformally identical to their
root cells, and are crossed in series.  Apply the extremal-length Series Law.
\end{proof}

\begin{remark}
No block-independent lower bound for $R_*$ is needed below.  The physical
materialization depth adapts to the actual positive value.
\end{remark}

\section{Adaptive nodal response}
Let $S=\{\omega_1,\dots,\omega_s\}$ be the end addresses read by a finite proof
horizon.  Choose a depth beyond their common prefixes and, for address $i$, an
additional serial tail of length $\ell_i$.  Cut across the parent-side end of
each disjoint tail.  Let $\Lambda_{\rm phys}$ be the connected response and
$\Lambda_{\rm node}$ the block-diagonal response of the cut normalization.
Set
\[
 \mathcal R_i=\sum_{j=1}^{\ell_i}R_{i,j},\qquad c_i=\mathcal R_i^{-1}.
\]

\begin{theorem}[Quantitative nodal-response approximation]
On the quotient by componentwise constants,
\[
 0\preceqL\Lambda_{\rm phys}-\Lambda_{\rm node}.
\]
For every terminal vector $x$,
\[
 0\le x^T(\Lambda_{\rm phys}-\Lambda_{\rm node})x
 \le\left(\sum_{i=1}^sc_i\right)\osc(x)^2.
\]
Consequently,
\[
 \|\Lambda_{\rm phys}-\Lambda_{\rm node}\|_{\op}
 \le2\sum_{i=1}^sc_i
 \le2\sum_{i=1}^s\frac1{\ell_iR_{*,i}}.
\]
\end{theorem}

\begin{proof}
Cutting removes nonnegative coupling energy and permits independent constants
on the normalization components, proving the lower inequality.  Take the
energy minimizers on the cut components.  By the maximum principle, their free
attachment values lie between the global minimum and maximum terminal values.
Interpolate through neck $i$ by its two-terminal energy minimizer.  This costs
at most $c_i\osc(x)^2$.  Summing gives the quadratic-form bound.  Finally,
$\osc(x)^2\le2\|x\|_2^2$.
\end{proof}

\begin{corollary}[Explicit depth choice]
For a prescribed tolerance $\varepsilon>0$, it suffices to choose
\[
 \sum_i\frac{2}{\ell_iR_{*,i}}\le\varepsilon.
\]
If one lower bound $R_*$ applies to all selected tails, the uniform choice
\[
 \ell_i\ge\left\lceil\frac{2s}{\varepsilon R_*}\right\rceil
\]
gives
$\|\Lambda_{\rm phys}-\Lambda_{\rm node}\|_{\op}\le\varepsilon$.
\end{corollary}

\begin{remark}
The depth is finite side data selected after the finite proof horizon is
frozen and from a predeclared error budget.  The normalized end object itself
does not depend on this choice.
\end{remark}

\section{Finite-horizon and transition consequences}
Assume a finite normalized diagram has actual intrinsic branches, one
quadratic relative stage, critical-spine retention, a protected root packet and
port certificate, serially exhaustive necks, and a prescribed response error.

\begin{theorem}[Direct-decoration finite-horizon inflation]
The diagram has an actual finite physical realization with locally saturated
pseudo-quadratic-like stages, selected relative degree two, exact chronological
critical word and complete markings, a protected finite carrier, exact
fiberwise response modules, and total end-plumbing response error below the
prescribed tolerance.
\end{theorem}

\begin{proof}
Use the depth-graded realization theorem, the critical-spine telescope and
marked assembly, the finite protected-carrier theorem, exact fiberwise response,
and the adaptive depth corollary.
\end{proof}

Let $L$ be an error-stable homogeneous gauge satisfying
\[
 L(Z+\Delta)\le L(Z)+C_L\|\Delta\|.
\]
Assume the root block satisfies
\[
 L(Z_{\rm out})\le\mu L(Z_{\rm in})+E_{\rm root},
 \qquad 0<\mu<1.
\]

\begin{theorem}[Strict transition transfer]
For any preassigned $E_{\rm pl}>0$, choose the materialization depth so that
$\|\Delta_{\rm pl}\|\le E_{\rm pl}/C_L$.  Then
\[
 L(Z^{\rm vir}_{\rm out})
 \le\mu L(Z^{\rm vir}_{\rm in})+E_{\rm root}+E_{\rm pl}.
\]
Equivalently,
\[
 L_{\rm in}\ge\mu^{-1}L_{\rm out}
 -\mu^{-1}(E_{\rm root}+E_{\rm pl}).
\]
The homogeneous backward factor $\lambda=\mu^{-1}>1$ is unchanged by the
virtual address tail.
\end{theorem}

\begin{proof}
The normalized univalent tail is an identity chart.  The only added response
term is the plumbing perturbation, controlled by the preceding theorem.  Apply
error stability and rearrange.
\end{proof}

\subsection{Planar-free finite-interface execution}
Call a block finite-interface when every proof call reads a finite depth-graded
diagram; end operations are conformally natural; cross-fiber interaction is
read only through a finite protected carrier, finite-electrode response, and
retained finite Value features; inequalities are stable under a prescribed
response error; and the horizon and error budget are frozen before depth is
chosen.

\begin{theorem}[Pro-exact execution sufficiency]
A finite-interface direct-decoration block supplied with its actual root
branches, intrinsic finite atlases, critical spine, packet/port certificate,
and serial necks has, for every $E_{\rm pl}>0$, an actual finite-cylinder
execution with exact dynamical and marked data, one protected finite carrier,
and total cross-fiber response error at most $E_{\rm pl}$.  Every
error-stable strict block estimate survives with that additive charge.  No
point-valued landing map or nondegenerate planar infinite-address component is
used.
\end{theorem}

\begin{proof}
Unfold shift-grading conflicts and apply finite physical realization.  Use the
protected-carrier theorem, distribute the finite error budget among the
selected tails, and choose their depths by the nodal-response estimate.  Exact
fiber data and strict-transition transfer finish the execution.
\end{proof}

Thus planar virtual-component identification is not a prerequisite for any
module satisfying this finite-interface contract.  Universal source
authorization and a uniformly adequate observable interface remain open.

\section{Promotion boundary}
The following split is required to avoid theorem laundering:
\begin{center}
\begin{tabular}{p{0.34\linewidth}p{0.27\linewidth}p{0.27\linewidth}}
\toprule
Original node & Proved part & Still open \\
\midrule
Disconnected marking transfer & Rescaled end transfer and finite realization & Identification with the intended planar virtual component \\
Cantor package realization & Fibered pro-exact package & Actual planar Cantor family and global protected coupling \\
Prefix-local scale selector & Intrinsic finite atlas compiler & Universal source authorization at every mixed scale \\
\bottomrule
\end{tabular}
\end{center}

The present results also do not prove a universally bounded future-observable
response/\allowbreak WAD/\allowbreak Value dimension, a strict source alternative for every mixed
block, MBTT, coercive compactness, or mixed-tower rigidity.

\section{Next theorem}
The next source-facing theorem should universally authorize the finite
normalized interface for every mixed block: actual root portraits, intrinsic
finite atlases, critical spines, protected packet/port and serial-neck
certificates, and a proof-complete demand schema.  It should then construct a
universally controlled future-observable response/WAD/Value quotient of the
complete coupled interaction, with a strict root alternative and discounted
error budget.  The planar identification of infinite-address components
remains an open geometric refinement, but Theorem~8.3 shows that it is not a
logical prerequisite for this finite-interface route.  The remaining global
steps are MBTT, geometric coercivity, and mixed-tower rigidity.

\end{document}
